\documentclass[12pt]{article}
\usepackage[utf8]{inputenc}
\usepackage[T1]{fontenc}
\usepackage[margin=1in]{geometry}
\usepackage{setspace}
\usepackage{xcolor}
\usepackage[colorlinks=true, urlcolor=blue!60!black, linkcolor=blue!60!black]{hyperref}
\setlength{\parindent}{0em}
\setlength{\parskip}{0.6em}
\renewcommand{\baselinestretch}{1.0}
\begin{document}
\begin{center}
{\LARGE \textit{United States v. Lopez}, 514 U.S. 549 (1995)\par}
\vspace{0.5em}
{\large Case Brief --- Two-Page Edition\par}
\end{center}

\section{Memory Jogger}
For the first time in nearly six decades, the Court held a federal statute beyond the commerce power: possessing a gun in a school zone is not economic activity and was not tied by a jurisdictional element to interstate commerce (\href{https://www.courtlistener.com/opinion/117927/united-states-v-lopez/}{514 U.S. 549}).

\section{Detailed Case Facts}
Alfonso Lopez, Jr., a San Antonio high-school senior, brought a concealed handgun to school in 1992. Federal agents charged him under the Gun-Free School Zones Act of 1990, which barred knowingly possessing a firearm within 1,000 feet of a school and contained neither congressional findings about commerce nor a jurisdictional element.

\section{Procedural History}
The District Court upheld the Act and convicted Lopez after a bench trial; the Fifth Circuit reversed. The Supreme Court affirmed.

\section{Judicial Votes}
Affirmed, 5--4. Rehnquist wrote for O'Connor, Scalia, Kennedy, and Thomas; Kennedy (with O'Connor) and Thomas concurred; Stevens, Souter, and Breyer (with Stevens, Souter, and Ginsburg) dissented.

\section{Holding}
Congress may regulate the channels and instrumentalities of interstate commerce and activities that ``substantially affect'' it. A criminal ban on local, noneconomic conduct, with no jurisdictional element and only an attenuated link to commerce, exceeds that power.

\section{Analysis of Opinions}
Rehnquist found that the Act regulated no economic enterprise, was not part of a larger market scheme as in \textit{Wickard}, and lacked a jurisdictional hook; the Government's ``costs of crime'' and ``national productivity'' theories had no limiting principle. Kennedy called the holding ``necessary though limited,'' stressing political accountability and the States' role in education. Thomas urged returning to the original meaning of ``commerce.'' Souter warned of \textit{Lochner}-era formalism, and Breyer found a rational basis linking school violence to the national economy.

\section{Examples: Future Applications}
\textbf{Same-Side:} a federal crime of high-school truancy, justified by national productivity, fails. \textbf{Opposite-Side:} banning home marijuana cultivation within a comprehensive drug-market scheme is valid, as \textit{Gonzales v. Raich} held. \textbf{Fence-Sitter:} the 1996 version of the Act, which requires a gun that has moved in interstate commerce.

\section{Critique}
Randy Barnett defends \textit{Lopez} as restoring enumerated powers, though some originalists call it a half-measure that kept \textit{Wickard}. Erwin Chemerinsky and others argue it revived pre-1937 formalism without a workable economic/noneconomic line.

\section{Key Quotations}
\begin{quote}
``To uphold the Government's contentions here, we would have to pile inference upon inference in a manner that would bid fair to convert congressional authority under the Commerce Clause to a general police power of the sort retained by the States.''
\end{quote}
\noindent --- Chief Justice Rehnquist, Opinion of the Court.

\end{document}
